\chapterimage{onions.jpg}


\chapter{Una pila de capes}\label{sec:stack_layers}

La comunicació moderna entre ordinadors s'aconsegueix fent servir una
\conceptRef{stack}{pila} de \conceptRef{layer}{capes} de xarxa.
Cada pila se centra en alguns problemes i és responsable de 
proporcionar certes funcionalitats, \eg, 
l'adreçament i la \conceptRef{sequencing}{seqüenciació} correcta (presentats en els capítols anteriors).

Cada \conceptRef{layer}{capa} defineix una \textit{parella} de funcions \inlineCode{send}/\inlineCode{receive},
que és responsable d'algunes d'aquestes funcionalitats.
% 
Cada capa fa servir el \inlineCode{send} i el \inlineCode{receive} de la capa inferior,
\ie, les funcionalitats d'una capa es basen en les funcionalitats de totes les capes que té per sota.

Per coordinar les capes i permetre una comunicació virtual entre elles,
cada capa defineix un format de paquet diferent, que \conceptRef{encapsulation}{encapsula} 
un paquet de la capa superior.


\section{Enviament de dades}\label{sec:stack_layers:send}

Per implementar aquestes 
capacitats, cada capa transmet no només les dades (\inlineCode{data}) proporcionades per la capa superior, 
sinó també una capçalera amb tota la metainformació necessària per a aquella capacitat. Així, cada capa ha de definir
el format de paquet, afegint una certa sobrecàrrega (típicament una capçalera) al payload (càrrega útil), \eg:

\begin{center}
    \begin{bytefield}{36}
     \bitbox{12}{Header} & \bitbox{24}{Payload}
    \end{bytefield}
\end{center}



Quan un programa d'usuari vol enviar dades, fa servir el mètode \inlineCode{send} de la capa \textit{superior} de la pila,
cosa que acaba fent que la capa física transmeti dades, \eg, a través d'un cable o d'una antena.
La figura següent mostra l'estructura d'una \conceptRef{stack}{pila} de 4 capes:

\begin{center}
\includegraphics[width=\linewidth]{encapsulation_send.pdf}
\end{center}


Els paquets que produeix cada capa s'anomenen genèricament \concept{PDU}s (Protocol Data Unit, unitat de dades del protocol).
Les PDUs de cada capa tenen noms específics com ara \concept{frame} (trama), \conceptRef{datagram}{datagrama} o segment.
Aquests noms s'han de fer servir amb precisió en contextos tècnics.

Cada capa inclou la \concept{PDU} de la capa superior dins la definició del format de paquet de la seva PDU.
Aquest procés s'anomena \conceptRef{encapsulation}{encapsulació}.

\begin{exercise}

El programa d'usuari vol enviar $420$~bytes de dades fent servir la \conceptRef{stack}{pila} de la figura anterior.
Suposa una mida de capçalera per a les capes 2, 3 i 4 de $14$~bytes, $20$~bytes i $20$~bytes, respectivament.
\begin{itemize}
    \item Quina és la mida del payload i la mida de la sobrecàrrega de la \concept{PDU} de cada \conceptRef{layer}{capa}?
    \item Quina és la mida del payload i la mida de la sobrecàrrega de tota la \conceptRef{stack}{pila}?
\end{itemize}
\end{exercise}

\section{Recepció de dades}\label{sec:stack_layers:receive}

Quan un paquet arriba a la \conceptRef{stack}{pila} de xarxa del destinatari, es processa en ordre invers,
\ie, les capes es recorren de baix a dalt. Cada capa inspecciona \textit{només} la capçalera corresponent a aquella capa.
Si tot és correcte, la capa \conceptRef{decapsulate}{elimina la capçalera} i passa la resta del paquet a la capa superior.
Al seu torn, la capa de més amunt passa les dades restants (el payload) al destinatari.
Continuant amb l'exemple de la pila de 4 capes:

\begin{center}
\includegraphics[width=\linewidth]{encapsulation_receive.pdf}
\end{center}


En una \conceptRef{stack}{pila de xarxa}, només la capa inferior interactua directament amb l'exterior,
\eg, només la targeta de xarxa física d'un ordinador està connectada a un cable que porta dades cap a altres ordinadors.
% 
Alhora, les dades enviades per la capa inferior inclouen les \conceptRef{PDU}{PDUs} de totes les capes
mitjançant l'\conceptRef{encapsulation}{encapsulació}.
% 
Gràcies a això, hi ha una comunicació virtual entre cada capa d'un extrem i la seva homòloga a l'altre
extrem de la comunicació. La figura següent representa la comunicació real (línia contínua) i virtual (línia discontínua) 
del nostre exemple:

\begin{center}
    \includegraphics[width=0.5\linewidth]{horizontal_layer_communication.pdf}
\end{center}


\begin{exercise}
Un altre exemple de \conceptRef{stack}{pila} de xarxa té 3 capes: L1, L2 i L3, de baix a dalt.

La \concept{PDU} de la capa L2 és la següent:

\begin{center}
\begin{bytefield}{32}
\bitheader{0,7,8,15,16,23,24,31}\\
\bitbox{16}{Destination address} & \bitbox{16}{Payload length (bytes)} \\
\wordbox{2}{{Payload (variable length)\\$\vdots$}}
\end{bytefield}
\end{center}

La \concept{PDU} de la capa L3 és la següent:

\begin{center}
\begin{bytefield}{32}
\bitheader{0,7,8,15,16,23,24,31}\\
\bitbox{16}{Fragment offset} & \bitbox{16}{Payload length (bytes)} \\
\wordbox{2}{{Payload (variable length)\\$\vdots$}}
\end{bytefield}
\end{center}

La capa L1 rep un paquet de dades i la seva funció \inlineCode{receive} retorna les dades següents
(suposa enters sense signe i ordre big endian (primer el byte més significatiu) quan calgui:
\begin{center}
\otherBase{01 23 00 0a 01 00 00 06 00 aa 00 bb 43 20}
\end{center}
\begin{itemize}
\item Quin payload de dades s'ha enviat al programa destinatari en aquest paquet?\\[-0.25cm]

\item Si aquest és l'últim paquet de la comunicació, quantes dades (payload) s'han enviat
  al programa d'usuari destinatari?\\[-0.25cm]
  
\item Com a màxim, quants ordinadors diferents es poden connectar fent servir aquesta \conceptRef{stack}{pila}?\\[-0.25cm]

\item Enviem un fitxer com a payload d'aquesta pila. Quina és la seva mida màxima possible?
\end{itemize}
\label{ex:how_it_is_done:three_layer_network}
\end{exercise}

\begin{center}
\showCode{snippets/encapsulation.py}
\end{center}

\begin{exercise}
El codi anterior agafa unes dades i escriu els bytes que s'enviarien a la capa 1 
(la \conceptRef{physical layer}{capa física}) de la \conceptRef{stack}{pila de xarxa} de l'exercici~\ref{ex:how_it_is_done:three_layer_network}.
Implementa l'extrem receptor de la pila de la manera següent: 

\begin{itemize}
\item Implementa una funció \inlineCode{layer2_decapsulate(data: bytes) -> (int, bytes)} que rebi una PDU de capa 2
  i retorni la tupla \inlineCode{(address, payload)}, on \inlineCode{address} és l'adreça de destinació i 
  \inlineCode{payload} és la PDU de capa 3 encapsulada.\\[-0.25cm]
  
\item Implementa una funció \inlineCode{layer3_decapsulate(data: bytes) -> (int, bytes)} que rebi una PDU de capa 3
  i retorni la tupla \inlineCode{(offset, payload)}, on \inlineCode{offset} és l'offset (desplaçament) del payload, i 
  \inlineCode{payload} és el fragment de dades enviat per l'altre extrem de la comunicació.\\[-0.25cm]
  
\item Implementa una funció \inlineCode{stack_receive(data: bytes)} que descodifiqui una PDU de capa 1 (fent servir les dues funcions anteriors)
  i imprimeixi una línia de missatge semblant a 
  \begin{center}
  \inlineCode{"Received a chunk of data. Length: 12 bytes, Address: 0xabcd, Offset: 0."}
  \end{center}\ \\[-0.75cm]
  
\item Completa la funció \inlineCode{stack_send} perquè les dades d'entrada es divideixin en fragments de com a màxim 100 bytes, 
      es produeixi una PDU de capa 1 vàlida per a cadascun d'aquests fragments,
      i aquestes PDUs de capa 1 es passin a \inlineCode{stack_receive}.\\[-0.25cm]
      
\item La teva implementació funcionaria si les PDUs de capa 1 s'ordenessin aleatòriament abans de passar-les a la pila receptora?
\end{itemize}
% 
\label{ex:how_it_is_done:implementation_three_layer_network}
\end{exercise}

\begin{remark}
La biblioteca \inlineCode{struct} de Python és útil, però no obligatòria, per donar format als bytes dels paquets.
\end{remark}

